\documentclass{article}
\usepackage{ctex}
\begin{document}

本系列手册主要涵盖最新的英特尔\textsuperscript{\textregistered}  64 位与 IA-32 架构处理器相关信息，包括：

\begin{itemize}
\item 奔腾\textsuperscript{\textregistered}  处理器
\item P6 系列处理器
\item 奔腾\textsuperscript{\textregistered}  4 处理器
\item 奔腾\textsuperscript{\textregistered}  M 处理器
\item 英特尔\textsuperscript{\textregistered}  至强\textsuperscript{\textregistered}  处理器
\item 奔腾\textsuperscript{\textregistered}  D 处理器
\item 奔腾\textsuperscript{\textregistered}  处理器至尊版
\item 64 位英特尔\textsuperscript{\textregistered}  至强\textsuperscript{\textregistered}  处理器
\item 英特尔\textsuperscript{\textregistered}  酷睿$^{\rm{TM}}$ 双核处理器
\item 英特尔\textsuperscript{\textregistered}  酷睿$^{\rm{TM}}$ 单核处理器
\item 双核英特尔\textsuperscript{\textregistered}  至强\textsuperscript{\textregistered}  处理器 LV
\item 英特尔\textsuperscript{\textregistered}  酷睿$^{\rm{TM}}$ 2 双核处理器
\item 英特尔\textsuperscript{\textregistered}  酷睿$^{\rm{TM}}$ 2 四核处理器 Q6000 系列
\item 英特尔\textsuperscript{\textregistered}  至强\textsuperscript{\textregistered}  处理器 3000、3200 系列
\item 英特尔\textsuperscript{\textregistered}  至强\textsuperscript{\textregistered}  处理器 5000 系列
\item 英特尔\textsuperscript{\textregistered}  至强\textsuperscript{\textregistered}  处理器 5100、5300 系列
\item 英特尔\textsuperscript{\textregistered}  酷睿$^{\rm{TM}}$ 2 至尊版处理器 X7000 和 X6800 系列
\item 英特尔\textsuperscript{\textregistered}  酷睿$^{\rm{TM}}$ 2 至尊版处理器 QX6000 系列
\item 英特尔\textsuperscript{\textregistered}  至强\textsuperscript{\textregistered}  处理器 7100 系列
\item 英特尔\textsuperscript{\textregistered}  奔腾\textsuperscript{\textregistered}  双核处理器
\item 英特尔\textsuperscript{\textregistered}  至强\textsuperscript{\textregistered}  处理器 7200、7300 系列
\item 英特尔\textsuperscript{\textregistered}  至强\textsuperscript{\textregistered}  处理器 5200、5400、7400 系列
\item 英特尔\textsuperscript{\textregistered}  酷睿$^{\rm{TM}}$ 2 至尊版处理器 QX9000 和 X9000 系列
\item 英特尔\textsuperscript{\textregistered}  酷睿$^{\rm{TM}}$ 2 四核处理器 Q9000 系列
\item 英特尔\textsuperscript{\textregistered}  酷睿$^{\rm{TM}}$ 2 双核处理器 E8000、T9000 系列
\item 英特尔\textsuperscript{\textregistered}  凌动$^{\rm{TM}}$ 处理器家族
\item 英特尔\textsuperscript{\textregistered}  凌动$^{\rm{TM}}$ 处理器 200、300、D400、D500、D2000、N200、N400、N2000、E2000、Z500、Z600、Z2000、C1000 系列基于 45 纳米和 32 纳米工艺制造
\item 英特尔\textsuperscript{\textregistered}  酷睿$^{\rm{TM}}$ i7 处理器
\item 英特尔\textsuperscript{\textregistered}  酷睿$^{\rm{TM}}$ i5 处理器
\item 英特尔\textsuperscript{\textregistered}  至强\textsuperscript{\textregistered}  处理器 E7-8800/4800/2800 产品家族
\item 英特尔\textsuperscript{\textregistered}  酷睿$^{\rm{TM}}$ i7-3930K 处理器
\item 第二代英特尔\textsuperscript{\textregistered}  酷睿$^{\rm{TM}}$ i7-2xxx、英特尔\textsuperscript{\textregistered}  酷睿$^{\rm{TM}}$ i5-2xxx、英特尔\textsuperscript{\textregistered}  酷睿$^{\rm{TM}}$ i3-2xxx 处理器系列
\item 英特尔\textsuperscript{\textregistered}  至强\textsuperscript{\textregistered}  处理器 E3-1200 产品家族
\item 英特尔\textsuperscript{\textregistered}  至强\textsuperscript{\textregistered}  处理器 E5-2400/1400 产品家族
\item 英特尔\textsuperscript{\textregistered}  至强\textsuperscript{\textregistered}  处理器 E5-4600/2600/1600 产品家族
\item 第三代英特尔\textsuperscript{\textregistered}  酷睿$^{\rm{TM}}$ 处理器
\item 英特尔\textsuperscript{\textregistered}  至强\textsuperscript{\textregistered}  处理器 E3-1200 v2 产品家族
\item 英特尔\textsuperscript{\textregistered}  至强\textsuperscript{\textregistered}  处理器 E5-2400/1400 v2 产品家族
\item 英特尔\textsuperscript{\textregistered}  至强\textsuperscript{\textregistered}  处理器 E5-4600/2600/1600 v2 产品家族
\item 英特尔\textsuperscript{\textregistered}  至强\textsuperscript{\textregistered}  处理器 E7-8800/4800/2800 v2 产品家族
\item 第四代英特尔\textsuperscript{\textregistered}  酷睿$^{\rm{TM}}$ 处理器
\item 英特尔\textsuperscript{\textregistered}  酷睿$^{\rm{TM}}$ M 处理器家族
\item 英特尔\textsuperscript{\textregistered}  酷睿$^{\rm{TM}}$ i7-59xx 处理器至尊版
\item 英特尔\textsuperscript{\textregistered}  酷睿$^{\rm{TM}}$ i7-49xx 处理器至尊版
\item 英特尔\textsuperscript{\textregistered}  至强\textsuperscript{\textregistered}  处理器 E3-1200 v3 产品家族
\item 英特尔\textsuperscript{\textregistered}  至强\textsuperscript{\textregistered}  处理器 E5-2600/1600 v3 产品家族
\item 第五代英特尔\textsuperscript{\textregistered}  酷睿$^{\rm{TM}}$ 处理器
\item 英特尔\textsuperscript{\textregistered}  至强\textsuperscript{\textregistered}  处理器 D-1500 产品家族
\item 英特尔\textsuperscript{\textregistered}  至强\textsuperscript{\textregistered}  处理器 E5 v4 家族
\item 英特尔\textsuperscript{\textregistered}  凌动$^{\rm{TM}}$ 处理器 X7-Z8000 和 X5-Z8000 系列
\item 英特尔\textsuperscript{\textregistered}  凌动$^{\rm{TM}}$ 处理器 Z3400 系列
\item 英特尔\textsuperscript{\textregistered}  凌动$^{\rm{TM}}$ 处理器 Z3500 系列
\item 第六代英特尔\textsuperscript{\textregistered}  酷睿$^{\rm{TM}}$ 处理器
\item 英特尔\textsuperscript{\textregistered}  至强\textsuperscript{\textregistered}  处理器 E3-1500m v5 产品家族
\item 第七代英特尔\textsuperscript{\textregistered}  酷睿$^{\rm{TM}}$ 处理器
\item 英特尔\textsuperscript{\textregistered}  至强融核$^{\rm{TM}}$ 处理器 3200、5200、7200 系列
\item 英特尔\textsuperscript{\textregistered}  至强\textsuperscript{\textregistered}  可扩展处理器家族
\item 第八代英特尔\textsuperscript{\textregistered}  酷睿$^{\rm{TM}}$ 处理器
\item 英特尔\textsuperscript{\textregistered}  至强融核$^{\rm{TM}}$ 处理器 7215、7285、7295 系列
\item 英特尔\textsuperscript{\textregistered}  至强\textsuperscript{\textregistered}  E 处理器
\item 第九代英特尔\textsuperscript{\textregistered}  酷睿$^{\rm{TM}}$ 处理器
\item 第二代英特尔\textsuperscript{\textregistered}  至强\textsuperscript{\textregistered}  可扩展处理器家族
\item 第十代英特尔\textsuperscript{\textregistered}  酷睿$^{\rm{TM}}$ 处理器
\item 第十一代英特尔\textsuperscript{\textregistered}  酷睿$^{\rm{TM}}$ 处理器
\item 第三代英特尔\textsuperscript{\textregistered}  至强\textsuperscript{\textregistered}  可扩展处理器家族
\item 第十二代英特尔\textsuperscript{\textregistered}  酷睿$^{\rm{TM}}$ 处理器
\item 第十三代英特尔\textsuperscript{\textregistered}  酷睿$^{\rm{TM}}$ 处理器
\item 第四代英特尔\textsuperscript{\textregistered}  至强\textsuperscript{\textregistered}  可扩展处理器家族
\item 第五代英特尔\textsuperscript{\textregistered}  至强\textsuperscript{\textregistered}  可扩展处理器家族
\item 英特尔\textsuperscript{\textregistered}  酷睿$^{\rm{TM}}$ Ultra 7 处理器
\item 英特尔\textsuperscript{\textregistered}  至强\textsuperscript{\textregistered}  6 E 核处理器
\item 英特尔\textsuperscript{\textregistered}  至强\textsuperscript{\textregistered}  6 P 核处理器
\item 英特尔\textsuperscript{\textregistered}  系列 2 酷睿$^{\rm{TM}}$ Ultra 处理器
\end{itemize}

P6 系列处理器是基于 P6 家族微架构的 IA-32 处理器。这包括奔腾\textsuperscript{\textregistered} Pro、奔腾\textsuperscript{\textregistered} II、奔腾\textsuperscript{\textregistered} III 和奔腾\textsuperscript{\textregistered} III 至强\textsuperscript{\textregistered} 处理器。

奔腾\textsuperscript{\textregistered} 4、奔腾\textsuperscript{\textregistered} D 和奔腾\textsuperscript{\textregistered} 处理器至尊版基于英特尔 NetBurst\textsuperscript{\textregistered} 微架构。

大多数早期英特尔\textsuperscript{\textregistered} 至强\textsuperscript{\textregistered} 处理器基于英特尔 NetBurst\textsuperscript{\textregistered} 微架构。英特尔至强处理器 5000、7100 系列基于英特尔 NetBurst\textsuperscript{\textregistered} 微架构。

英特尔\textsuperscript{\textregistered} 酷睿$^{\rm{TM}}$ 双核、英特尔\textsuperscript{\textregistered} 酷睿$^{\rm{TM}}$ 单核和双核英特尔\textsuperscript{\textregistered} 至强\textsuperscript{\textregistered} 处理器 LV 基于改进的奔腾\textsuperscript{\textregistered} M 处理器微架构。

英特尔\textsuperscript{\textregistered} 至强\textsuperscript{\textregistered} 处理器 3000、3200、5100、5300、7200 和 7300 系列，英特尔\textsuperscript{\textregistered} 奔腾\textsuperscript{\textregistered} 双核，英特尔\textsuperscript{\textregistered} 酷睿$^{\rm{TM}}$ 2 双核，英特尔\textsuperscript{\textregistered} 酷睿$^{\rm{TM}}$ 2 四核和英特尔\textsuperscript{\textregistered} 酷睿$^{\rm{TM}}$ 2 至尊版处理器基于英特尔\textsuperscript{\textregistered} 酷睿$^{\rm{TM}}$ 微架构。

英特尔\textsuperscript{\textregistered} 至强\textsuperscript{\textregistered} 处理器 5200、5400、7400 系列，英特尔\textsuperscript{\textregistered} 酷睿$^{\rm{TM}}$ 2 四核处理器 Q9000 系列，以及英特尔\textsuperscript{\textregistered} 酷睿$^{\rm{TM}}$ 2 至尊版处理器 QX9000、X9000 系列，英特尔\textsuperscript{\textregistered} 酷睿$^{\rm{TM}}$ 2 处理器 E8000 系列基于增强型英特尔\textsuperscript{\textregistered} 酷睿$^{\rm{TM}}$ 微架构。

英特尔\textsuperscript{\textregistered} 凌动$^{\rm{TM}}$ 处理器 200、300、D400、D500、D2000、N200、N400、N2000、E2000、Z500、Z600、Z2000、C1000 系列基于英特尔凌动$^{\rm{TM}}$ 微架构并支持英特尔 64 位架构。

P6 家族、奔腾\textsuperscript{\textregistered} M、英特尔\textsuperscript{\textregistered} 酷睿$^{\rm{TM}}$ 单核、英特尔\textsuperscript{\textregistered} 酷睿$^{\rm{TM}}$ 双核处理器、双核英特尔\textsuperscript{\textregistered} 至强\textsuperscript{\textregistered} 处理器 LV，以及早期世代的奔腾 4 和英特尔至强处理器支持 IA-32 架构。英特尔\textsuperscript{\textregistered} 凌动$^{\rm{TM}}$ 处理器 Z5xx 系列支持 IA-32 架构。

英特尔\textsuperscript{\textregistered} 至强\textsuperscript{\textregistered} 处理器 3000、3200、5000、5100、5200、5300、5400、7100、7200、7300、7400 系列，英特尔\textsuperscript{\textregistered} 酷睿$^{\rm{TM}}$ 2 双核、英特尔\textsuperscript{\textregistered} 酷睿$^{\rm{TM}}$ 2 至尊版、英特尔\textsuperscript{\textregistered} 酷睿$^{\rm{TM}}$ 2 四核处理器，奔腾\textsuperscript{\textregistered} D 处理器，奔腾\textsuperscript{\textregistered} 双核处理器，以及更新世代的奔腾 4 和英特尔至强处理器家族支持英特尔\textsuperscript{\textregistered} 64 位架构。

英特尔\textsuperscript{\textregistered} 酷睿$^{\rm{TM}}$ i7 处理器和英特尔\textsuperscript{\textregistered} 至强\textsuperscript{\textregistered} 处理器 3400、5500、7500 系列基于 45 纳米 Nehalem 微架构。Westmere 微架构是 Nehalem 微架构的 32 纳米版本。英特尔\textsuperscript{\textregistered} 至强\textsuperscript{\textregistered} 处理器 5600 系列、英特尔至强处理器 E7 以及各种英特尔酷睿 i7、i5、i3 处理器基于 Westmere 微架构。这些处理器支持英特尔 64 位架构。

英特尔\textsuperscript{\textregistered} 至强\textsuperscript{\textregistered} 处理器 E5 家族、英特尔\textsuperscript{\textregistered} 至强\textsuperscript{\textregistered} 处理器 E3-1200 家族、英特尔\textsuperscript{\textregistered} 至强\textsuperscript{\textregistered} 处理器 E7-8800/4800/2800 产品家族、英特尔\textsuperscript{\textregistered} 酷睿$^{\rm{TM}}$ i7-3930K 处理器，以及第二代英特尔\textsuperscript{\textregistered} 酷睿$^{\rm{TM}}$ i7-2xxx、英特尔\textsuperscript{\textregistered} 酷睿$^{\rm{TM}}$ i5-2xxx、英特尔\textsuperscript{\textregistered} 酷睿$^{\rm{TM}}$ i3-2xxx 处理器系列基于 Sandy Bridge 微架构并支持英特尔 64 位架构。

英特尔\textsuperscript{\textregistered} 至强\textsuperscript{\textregistered} 处理器 E7-8800/4800/2800 v2 产品家族、英特尔\textsuperscript{\textregistered} 至强\textsuperscript{\textregistered} 处理器 E3-1200 v2 产品家族和第三代英特尔\textsuperscript{\textregistered} 酷睿$^{\rm{TM}}$ 处理器基于 Ivy Bridge 微架构并支持英特尔 64 位架构。

英特尔\textsuperscript{\textregistered} 至强\textsuperscript{\textregistered} 处理器 E5-4600/2600/1600 v2 产品家族、英特尔\textsuperscript{\textregistered} 至强\textsuperscript{\textregistered} 处理器 E5-2400/1400 v2 产品家族和英特尔\textsuperscript{\textregistered} 酷睿$^{\rm{TM}}$ i7-49xx 处理器至尊版基于 Ivy Bridge-E 微架构并支持英特尔 64 位架构。

英特尔\textsuperscript{\textregistered} 至强\textsuperscript{\textregistered} 处理器 E3-1200 v3 产品家族和第四代英特尔\textsuperscript{\textregistered} 酷睿$^{\rm{TM}}$ 处理器基于 Haswell 微架构并支持英特尔 64 位架构。

英特尔\textsuperscript{\textregistered} 至强\textsuperscript{\textregistered} 处理器 E5-2600/1600 v3 产品家族和英特尔\textsuperscript{\textregistered} 酷睿$^{\rm{TM}}$ i7-59xx 处理器至尊版基于 Haswell-E 微架构并支持英特尔 64 位架构。

英特尔\textsuperscript{\textregistered} 凌动$^{\rm{TM}}$ 处理器 Z8000 系列基于 Airmont 微架构。

英特尔\textsuperscript{\textregistered} 凌动$^{\rm{TM}}$ 处理器 Z3400 系列和英特尔\textsuperscript{\textregistered} 凌动$^{\rm{TM}}$ 处理器 Z3500 系列基于 Silvermont 微架构。

英特尔\textsuperscript{\textregistered} 酷睿$^{\rm{TM}}$ M 处理器家族、第五代英特尔\textsuperscript{\textregistered} 酷睿$^{\rm{TM}}$ 处理器、英特尔\textsuperscript{\textregistered} 至强\textsuperscript{\textregistered} 处理器 D-1500 产品家族和英特尔\textsuperscript{\textregistered} 至强\textsuperscript{\textregistered} 处理器 E5 v4 家族基于 Broadwell 微架构并支持英特尔 64 位架构。

英特尔\textsuperscript{\textregistered} 至强\textsuperscript{\textregistered} 可扩展处理器家族、英特尔\textsuperscript{\textregistered} 至强\textsuperscript{\textregistered} 处理器 E3-1500m v5 产品家族和第六代英特尔\textsuperscript{\textregistered} 酷睿$^{\rm{TM}}$ 处理器基于 Skylake 微架构并支持英特尔 64 位架构。

第七代英特尔\textsuperscript{\textregistered} 酷睿$^{\rm{TM}}$ 处理器基于 Kaby Lake 微架构并支持英特尔 64 位架构。

英特尔\textsuperscript{\textregistered} 凌动$^{\rm{TM}}$ 处理器 C 系列、英特尔\textsuperscript{\textregistered} 凌动$^{\rm{TM}}$ 处理器 X 系列、英特尔\textsuperscript{\textregistered} 奔腾\textsuperscript{\textregistered} 处理器 J 系列、英特尔\textsuperscript{\textregistered} 赛扬\textsuperscript{\textregistered} 处理器 J 系列和英特尔\textsuperscript{\textregistered} 赛扬\textsuperscript{\textregistered} 处理器 N 系列基于 Goldmont 微架构。

英特尔\textsuperscript{\textregistered} 至强融核$^{\rm{TM}}$ 处理器 3200、5200、7200 系列基于 Knights Landing 微架构并支持英特尔 64 位架构。

英特尔\textsuperscript{\textregistered} 奔腾\textsuperscript{\textregistered} 银牌处理器系列、英特尔\textsuperscript{\textregistered} 赛扬\textsuperscript{\textregistered} 处理器 J 系列和英特尔\textsuperscript{\textregistered} 赛扬\textsuperscript{\textregistered} 处理器 N 系列基于 Goldmont Plus 微架构。

第八代英特尔\textsuperscript{\textregistered} 酷睿$^{\rm{TM}}$ 处理器、第九代英特尔\textsuperscript{\textregistered} 酷睿$^{\rm{TM}}$ 处理器和英特尔\textsuperscript{\textregistered} 至强\textsuperscript{\textregistered} E 处理器基于 Coffee Lake 微架构并支持英特尔 64 位架构。

英特尔\textsuperscript{\textregistered} 至强融核$^{\rm{TM}}$ 处理器 7215、7285、7295 系列基于 Knights Mill 微架构并支持英特尔 64 位架构。

第二代英特尔\textsuperscript{\textregistered} 至强\textsuperscript{\textregistered} 可扩展处理器家族基于 Cascade Lake 产品并支持英特尔 64 位架构。

部分第十代英特尔\textsuperscript{\textregistered} 酷睿$^{\rm{TM}}$ 处理器基于 Ice Lake 微架构，部分基于 Comet Lake 微架构；两者均支持英特尔 64 位架构。

部分第十一代英特尔\textsuperscript{\textregistered} 酷睿$^{\rm{TM}}$ 处理器基于 Tiger Lake 微架构，部分基于 Rocket Lake 微架构；两者均支持英特尔 64 位架构。

部分第三代英特尔\textsuperscript{\textregistered} 至强\textsuperscript{\textregistered} 可扩展处理器家族处理器基于 Cooper Lake 产品，部分基于 Ice Lake 微架构；两者均支持英特尔 64 位架构。

第十二代英特尔\textsuperscript{\textregistered} 酷睿$^{\rm{TM}}$ 处理器基于 Alder Lake 性能混合架构并支持英特尔 64 位架构。

第十三代英特尔\textsuperscript{\textregistered} 酷睿$^{\rm{TM}}$ 处理器基于 Raptor Lake 性能混合架构并支持英特尔 64 位架构。

第四代英特尔\textsuperscript{\textregistered} 至强\textsuperscript{\textregistered} 可扩展处理器家族基于 Sapphire Rapids 微架构并支持英特尔 64 位架构。

第五代英特尔\textsuperscript{\textregistered} 至强\textsuperscript{\textregistered} 可扩展处理器家族基于 Emerald Rapids 微架构并支持英特尔 64 位架构。

英特尔\textsuperscript{\textregistered} 酷睿$^{\rm{TM}}$ Ultra 7 处理器基于 Meteor Lake 性能混合架构并支持英特尔 64 位架构。

英特尔\textsuperscript{\textregistered} 至强\textsuperscript{\textregistered} 6 E 核处理器基于 Sierra Forest 微架构并支持英特尔 64 位架构。

英特尔\textsuperscript{\textregistered} 至强\textsuperscript{\textregistered} 6 P 核处理器基于 Granite Rapids 微架构并支持英特尔 64 位架构。

英特尔\textsuperscript{\textregistered} 系列 2 酷睿$^{\rm{TM}}$ Ultra 处理器基于 Lunar Lake 性能混合架构并支持英特尔 64 位架构。

IA-32 架构是英特尔 32 位微处理器的指令集架构和编程环境。

英特尔\textsuperscript{\textregistered} 64 位架构是英特尔 32 位和 64 位架构的超集，是兼容 IA-32 架构的指令集架构和编程环境。


\end{document}